% GENERATED FILE. Edit canonical lessons, then run npm run generate:books.
% canonical-source-hash: fnv1a64:de8157eaabed1b0a
% canonical-lessons: BN-C452-ohongkar, BN-C452-demak, BN-C452-atmosomman, BN-C452-jahir-kora, BN-C452-opodostho-hooya

\chapter{Pride and Humiliation}
\label{ch:bn-pride-and-humiliation}
\hlchaptermodality{452}

\begin{chapteropening}
\textbf{By the end of this chapter:} \emph{I can talk about arrogance, boasting and being humiliated.}

Complete the last lesson of chapter 452: I can talk about arrogance, boasting and being humiliated.
\end{chapteropening}

\section[\texorpdfstring{\bn{অহংকার} (ôhôngkār)}{অহংকার (ôhôngkār)}]{\bn{অহংকার} (ôhôngkār) — arrogance; vanity}

\label{lesson:BN-C452-ohongkar}



\begin{quote}
Before the new one: say the Bengali for to hate, then the Bengali for to turn up one's nose.
\end{quote}

\subsection*{You'll want to know: \bn{অহংকার}}
\textbf{\bn{অহংকার}} — \emph{ôhôngkār} — \textquotedblleft{}arrogance; vanity\textquotedblright{}.

\begin{quote}
\bn{টাকা} \bn{হওয়ার} \bn{পরে} \bn{লোকটার} \bn{অহংকার} \bn{বাড়ল।}
\end{quote}

\begin{quote}
\emph{ṭākā hôoyār pôre lokṭār ôhôngkār bāṛlô.}
\end{quote}

\begin{quote}
Once he had money, the man's arrogance grew.
\end{quote}

\bn{অহংকারী}, arrogant. \bn{গর্ব} (taught) is a pride that can be good; \bn{অহংকার} never is.

\begin{cousinweb}
\textbf{\bn{অহংকার}} (\emph{ôhôngkār}) is a tatsama, Sanskrit \emph{ahaṃkāra}, \textquotedblleft{}I-making\textquotedblright{}: the sense of self, and so egotism (Monier-Williams).
\end{cousinweb}

\begin{grammarlens}[title={What you've built}]
One more word for describing what you lived and felt.
\end{grammarlens}

\subsection*{Your turn}
\begin{itemize}
\raggedright
  \item \emph{Say it:} \emph{ôhôngkār}
  \item \emph{Say it:} \emph{ôhôngkār}, once more
  \item \emph{Say it:} \emph{ṭākā hôoyār pôre lokṭār ôhôngkār bāṛlô.} — Once he had money, the man's arrogance grew.
  \item \emph{Recall:} say the Bengali for to hate, then the Bengali for to turn up one's nose, then say \emph{ôhôngkār} again
\end{itemize}

\begin{tcolorbox}[breakable,colback=teal!4,colframe=teal!35!black,title={Before you move on}]
What does \bn{অহংকার} mean? (\textquotedblleft{}Arrogance; vanity\textquotedblright{}.) Say it once more.
\end{tcolorbox}
\section[\texorpdfstring{\bn{দেমাক} (demāk)}{দেমাক (demāk)}]{\bn{দেমাক} (demāk) — conceit; airs}

\label{lesson:BN-C452-demak}



\begin{quote}
Before the new one: say the Bengali for to turn up one's nose, then the Bengali for arrogance.
\end{quote}

\subsection*{You'll want to know: \bn{দেমাক}}
\textbf{\bn{দেমাক}} — \emph{demāk} — \textquotedblleft{}conceit; airs\textquotedblright{}.

\begin{quote}
\bn{নতুন} \bn{গাড়ি} \bn{কিনে} \bn{প্রতিবেশীর} \bn{দেমাক} \bn{বাড়ল।}
\end{quote}

\begin{quote}
\emph{nôtun gāṛi kine prôtibeshīr demāk bāṛlô.}
\end{quote}

\begin{quote}
Having bought a new car, the neighbour began to put on airs.
\end{quote}

\begin{cousinweb}
\textbf{\bn{দেমাক}} (\emph{demāk}) is usually traced to Arabic \emph{dimāgh}, \textquotedblleft{}brain\textquotedblright{}, which took on the sense of pride and conceit in Persian.
\end{cousinweb}

\begin{grammarlens}[title={What you've built}]
One more word for describing what you lived and felt.
\end{grammarlens}

\subsection*{Your turn}
\begin{itemize}
\raggedright
  \item \emph{Say it:} \emph{demāk}
  \item \emph{Say it:} \emph{demāk}, once more
  \item \emph{Say it:} \emph{nôtun gāṛi kine prôtibeshīr demāk bāṛlô.} — Having bought a new car, the neighbour began to put on airs.
  \item \emph{Recall:} say the Bengali for to turn up one's nose, then the Bengali for arrogance, then say \emph{demāk} again
\end{itemize}

\begin{tcolorbox}[breakable,colback=teal!4,colframe=teal!35!black,title={Before you move on}]
What does \bn{দেমাক} mean? (\textquotedblleft{}Conceit; airs\textquotedblright{}.) Say it once more.
\end{tcolorbox}
\section[\texorpdfstring{\bn{আত্মসম্মান} (ātmôsômmān)}{আত্মসম্মান (ātmôsômmān)}]{\bn{আত্মসম্মান} (ātmôsômmān) — self-respect}

\label{lesson:BN-C452-atmosomman}



\begin{quote}
Before the new one: say the Bengali for arrogance, then the Bengali for conceit.
\end{quote}

\subsection*{You'll want to know: \bn{আত্মসম্মান}}
\textbf{\bn{আত্মসম্মান}} — \emph{ātmôsômmān} — \textquotedblleft{}self-respect\textquotedblright{}.

\begin{quote}
\bn{আত্মসম্মানের} \bn{জন্য} \bn{অরুণ} \bn{সাহায্য} \bn{নিল} \bn{না।}
\end{quote}

\begin{quote}
\emph{ātmôsômmāner jônyô ôrun sāhājjô nilô nā.}
\end{quote}

\begin{quote}
Out of self-respect Arun would not take help.
\end{quote}

\begin{cousinweb}
\textbf{\bn{আত্মসম্মান}} (\emph{ātmôsômmān}) puts Sanskrit \emph{ātman}, \textquotedblleft{}self\textquotedblright{} (Monier-Williams), in front of \bn{সম্মান} (taught), respect.
\end{cousinweb}

\begin{grammarlens}[title={What you've built}]
One more word for describing what you lived and felt.
\end{grammarlens}

\subsection*{Your turn}
\begin{itemize}
\raggedright
  \item \emph{Say it:} \emph{ātmôsômmān}
  \item \emph{Say it:} \emph{ātmôsômmān}, once more
  \item \emph{Say it:} \emph{ātmôsômmāner jônyô ôrun sāhājjô nilô nā.} — Out of self-respect Arun would not take help.
  \item \emph{Recall:} say the Bengali for arrogance, then the Bengali for conceit, then say \emph{ātmôsômmān} again
\end{itemize}

\begin{tcolorbox}[breakable,colback=teal!4,colframe=teal!35!black,title={Before you move on}]
What does \bn{আত্মসম্মান} mean? (\textquotedblleft{}Self-respect\textquotedblright{}.) Say it once more.
\end{tcolorbox}
\section[\texorpdfstring{\bn{জাহির} \bn{করা} (jāhir kôrā)}{জাহির করা (jāhir kôrā)}]{\bn{জাহির} \bn{করা} (jāhir kôrā) — to show off; to flaunt}

\label{lesson:BN-C452-jahir-kora}



\begin{quote}
Before the new one: say the Bengali for conceit, then the Bengali for self-respect.
\end{quote}

\subsection*{You'll want to know: \bn{জাহির} \bn{করা}}
\textbf{\bn{জাহির} \bn{করা}} — \emph{jāhir kôrā} — \textquotedblleft{}to show off; to flaunt\textquotedblright{}.

\begin{quote}
\bn{অরুণ} \bn{বন্ধুদের} \bn{সামনে} \bn{নতুন} \bn{ফোনটা} \bn{জাহির} \bn{করছিল।}
\end{quote}

\begin{quote}
\emph{ôrun bôndhuder sāmne nôtun phonṭā jāhir kôrchhilô.}
\end{quote}

\begin{quote}
Arun was showing off his new phone in front of his friends.
\end{quote}

Past (he/she …): \bn{তিনি} \bn{ছেলের} \bn{সাফল্য} \bn{জাহির} \bn{করলেন}, \emph{tini chheler sāphôlyô jāhir kôrlen}.

What is shown off is the object: \bn{জ্ঞান} \bn{জাহির} \bn{করা}, to show off what you know.

\begin{cousinweb}
\textbf{\bn{জাহির} \bn{করা}} (\emph{jāhir kôrā}) is \emph{jāhir}, \textquotedblleft{}made plain, on show\textquotedblright{}, usually traced to Arabic \emph{ẓāhir}, \textquotedblleft{}manifest\textquotedblright{}, through Persian, with \emph{kôrā}, \textquotedblleft{}to do\textquotedblright{} (CDIAL 2814): to put something where everyone must see it.
\end{cousinweb}

\begin{grammarlens}[title={What you've built}]
One more word for describing what you lived and felt.
\end{grammarlens}

\subsection*{Your turn}
\begin{itemize}
\raggedright
  \item \emph{Say it:} \emph{jāhir kôrā}
  \item \emph{Say it:} \emph{jāhir kôrā}, once more
  \item \emph{Say it:} \emph{ôrun bôndhuder sāmne nôtun phonṭā jāhir kôrchhilô.} — Arun was showing off his new phone in front of his friends.
  \item \emph{Recall:} say the Bengali for conceit, then the Bengali for self-respect, then say \emph{jāhir kôrā} again
\end{itemize}

\begin{tcolorbox}[breakable,colback=teal!4,colframe=teal!35!black,title={Before you move on}]
What does \bn{জাহির} \bn{করা} mean? (\textquotedblleft{}To show off; to flaunt\textquotedblright{}.) Say it once more.
\end{tcolorbox}
\section[\texorpdfstring{\bn{অপদস্থ} \bn{হওয়া} (ôpôdôsthô hôoyā)}{অপদস্থ হওয়া (ôpôdôsthô hôoyā)}]{\bn{অপদস্থ} \bn{হওয়া} (ôpôdôsthô hôoyā) — to be humiliated; to lose face}

\label{lesson:BN-C452-opodostho-hooya}



\begin{quote}
Before the new one: say the Bengali for self-respect, then the Bengali for to show off.
\end{quote}

\subsection*{You'll want to know: \bn{অপদস্থ} \bn{হওয়া}}
\textbf{\bn{অপদস্থ} \bn{হওয়া}} — \emph{ôpôdôsthô hôoyā} — \textquotedblleft{}to be humiliated; to lose face\textquotedblright{}.

\begin{quote}
\bn{বন্ধুদের} \bn{সামনে} \bn{বাবার} \bn{বকা} \bn{খেয়ে} \bn{অরুণ} \bn{অপদস্থ} \bn{হল।}
\end{quote}

\begin{quote}
\emph{bôndhuder sāmne bābār bôkā kheye ôrun ôpôdôsthô hôlô.}
\end{quote}

\begin{quote}
Scolded by his father in front of his friends, Arun felt humiliated.
\end{quote}

Past (he/she …): \bn{তিনি} \bn{সভায়} \bn{অপদস্থ} \bn{হলেন}, \emph{tini sôbhāy ôpôdôsthô hôlen}.

\begin{cousinweb}
\textbf{\bn{অপদস্থ} \bn{হওয়া}} (\emph{ôpôdôsthô hôoyā}) is Sanskrit \emph{apadastha}, \textquotedblleft{}not in its place, out of office\textquotedblright{} (Monier-Williams), which Bengali uses for someone put down in public, a sense Monier-Williams does not give, with \emph{hôoyā}, \textquotedblleft{}to be\textquotedblright{} (CDIAL 9416).
\end{cousinweb}

\begin{grammarlens}[title={What you've built}]
One more word for describing what you lived and felt.
\end{grammarlens}

\subsection*{Your turn}
\begin{itemize}
\raggedright
  \item \emph{Say it:} \emph{ôpôdôsthô hôoyā}
  \item \emph{Say it:} \emph{ôpôdôsthô hôoyā}, once more
  \item \emph{Say it:} \emph{bôndhuder sāmne bābār bôkā kheye ôrun ôpôdôsthô hôlô.} — Scolded by his father in front of his friends, Arun felt humiliated.
  \item \emph{Recall:} say the Bengali for self-respect, then the Bengali for to show off, then say \emph{ôpôdôsthô hôoyā} again
\end{itemize}

\begin{tcolorbox}[breakable,colback=teal!4,colframe=teal!35!black,title={Before you move on}]
What does \bn{অপদস্থ} \bn{হওয়া} mean? (\textquotedblleft{}To be humiliated; to lose face\textquotedblright{}.) Say it once more.
\end{tcolorbox}
